%!TEX root = ../main.tex

\section{Устойчивость положений равновесия}

\subsection{Упрощающие допущения}

Для описанной модели будем использовать разностную схему из работы \cite{ponomarev_stability}, где предварительно делается ряд допущений, упрощающих задачу. Кратко перечислим их.

Система \eqref{eq:Phi},~\eqref{eq:phi} рассматривается в замкнутой области $\clOmega = [0, W]_x \hm \times [0, H]_y \times I_z$, где $W, H > 0, \; I$~-- некоторый отрезок. Пусть $\epsilon_0(\vx) = \epsilon_0(x)$, а также задано начальное условие $\phi(\vx, 0) = \phi_0(x)$, то есть диэлектрическая проницаемость неповрежденной среды и начальное распределение фаз зависят только от $x$. На $\partial \Omega$ считаем заданным следующее граничное условие на $\phi$: $\phi|_{x = 0} = \phi_l(t)$, $\phi|_{x = W} = \phi_r(t)$, а также $\partflvn{\phi}= 0$ на <<гранях>> области $\clOmega$, перпендикулярных осям $y$ и $z$; следующее граничное условие~на~$\Phi$: $\Phi|_{y = 0} = \Phi^-$, $\Phi|_{y = H} = \Phi^+$, где $\Phi^-, \Phi^+ \in \Real$, $\Phi^+ \geqslant \Phi^-$, а также $\partflvn{\Phi} = 0$ на <<гранях>> $\clOmega$, перпендикулярных осям $x$ и $z$.

Учитывая описанные краевые условия, решение системы уравнений \eqref{eq:Phi},~\eqref{eq:phi} ищется в виде $\phi(\vx, t) = \phi(x, t)$, $\Phi(\vx, t) = \Phi(y, t)$. Решение $\Phi(\vx, t) = \Phi^- + (y / H)(\Phi^+ - \Phi^-)$ известно аналитически, и система, таким образом, сводится к единственному уравнению
\begin{equation}
	\cfrac{1}{m} \partt{\phi} = -F'(\phi; \norm{\nabla \Phi}) + \half \Gamma \partxx{\phi}
	\label{eq:one_dim}
\end{equation}
на функцию $\phi(x, t)$ в области $[0, W]_x \times [0, +\infty)_t$, с начальным условием
\begin{equation}
	\phi(x, 0) = \phi_0(x)
	\label{eq:one_dim_initial}{}
\end{equation}
и граничным условием
\begin{equation}
	\phi(0, t) = \phi_l(t), \quad \phi(W, t) = \phi_r(t).
	\label{eq:one_dim_marginal}
\end{equation}
Здесь $\norm{\nabla \Phi} = (\Phi^+ - \Phi^-) / H$ -- константа. В дальнейшем именно $\norm{\nabla \Phi}$ будет считаться параметром модели, конкретные значения $\Phi^+$ и $\Phi^-$ при этом неважны.

Так как в рассматриваемой постановке задачи $\norm{\nabla \Phi}$ -- константа, то вместо \linebreak $F(\phi; \norm{\nabla \Phi})$ будем писать $F(\phi)$.

Для простоты анализа везде далее $\epsilon_0$ считается константой.